\documentclass[11pt,a4paper]{article}
% Mise en page commune aux guides MyAnimLab (utilisateur, administrateur).
\usepackage[utf8]{inputenc}
\usepackage[T1]{fontenc}
\usepackage[french]{babel}
\usepackage[margin=2.2cm]{geometry}
\usepackage{enumitem,tabularx,booktabs}
\usepackage[table]{xcolor}
\usepackage{xurl}
\usepackage[hidelinks]{hyperref}
\usepackage{fancyhdr}
\usepackage[most]{tcolorbox}
\sloppy

\definecolor{accent}{HTML}{1F5F8B}
\definecolor{ok}{HTML}{2E7D32}
\definecolor{warn}{HTML}{C76A00}

% Chemin de menu : \menu{Animaux > Tournée quotidienne}
\newcommand{\menu}[1]{\textsf{\textbf{\color{accent}#1}}}
% Bouton ou libellé de l'interface
\newcommand{\btn}[1]{\textsf{«~#1~»}}
% Page ou fichier technique
\newcommand{\f}[1]{\nolinkurl{#1}}

% Procédure pas à pas
\newtcolorbox{procedure}[1]{colback=accent!4,colframe=accent,fonttitle=\bfseries\sffamily,
  title={#1},left=2mm,right=2mm,top=1mm,bottom=1mm,breakable}
% Règle à connaître
\newtcolorbox{regle}{colback=warn!6,colframe=warn,fonttitle=\bfseries\sffamily,
  title={À savoir},left=2mm,right=2mm,top=1mm,bottom=1mm,breakable}
% Traçabilité réglementaire
\newtcolorbox{trace}{colback=ok!6,colframe=ok,fonttitle=\bfseries\sffamily,
  title={Traçabilité},left=2mm,right=2mm,top=1mm,bottom=1mm,breakable}

\setlist[enumerate]{nosep,leftmargin=*}
\setlist[itemize]{nosep,leftmargin=*}

\pagestyle{fancy}
\fancyhf{}
\cfoot{\thepage}

\lhead{MyAnimLab --- Guide de l'administrateur}
\rhead{Version du 18 octobre 2026}

\title{\textbf{MyAnimLab}\\[4pt]\large Guide de l'administrateur}
\author{Administrateurs, responsables du bien-être animal (BEA)}
\date{18 octobre 2026}

\begin{document}
\maketitle
\thispagestyle{fancy}

\noindent Ce guide complète le \emph{Guide de l'utilisateur}. Il décrit
l'installation et la mise à jour, le paramétrage, les comptes, les
référentiels, la supervision quotidienne, les corrections réglementaires, le
volet BEA et la préparation d'une inspection.

\tableofcontents
\newpage

\section{Rôles}
\begin{center}
\rowcolors{2}{gray!8}{white}
\begin{tabularx}{\textwidth}{@{}lX@{}}
\toprule
\textbf{Rôle} & \textbf{Droits propres}\\
\midrule
Utilisateur & Toutes les tâches du guide utilisateur. Gestion des fichiers et retrait des cages auxquelles il est affecté.\\
Responsable BEA & En plus : analyse rétrospective BEA, bilan annuel, synthèse et validation des livrets de formation. Il reçoit les rappels de sevrage et d'actes en retard.\\
Administrateur & Tout, dont : comptes, paramètres du laboratoire, autorisations de projet, souches et lignées, corrections du registre, inventaires des stupéfiants, suppression définitive d'un animal, exports. Un administrateur est responsable BEA par défaut.\\
\bottomrule
\end{tabularx}
\end{center}

\section{Installation et mises à jour}

\subsection{Installation}
\begin{procedure}{Installation sur un serveur Ubuntu}
\begin{enumerate}
  \item Exécutez \f{setup/setup.sh}. Le script clone le dépôt, installe les
    dépendances, crée la base et importe le schéma avec le compte
    administrateur de MySQL. Il configure Apache (\texttt{AllowOverride
    All}) et donne au compte de l'application les seuls droits sur les
    données.
  \item Vérifiez que \f{https://votre-domaine/.env} répond \textbf{403 ou
    404}. Si la réponse est 200, le fichier \f{.htaccess} n'est pas pris en
    compte et vos identifiants sont exposés.
  \item Programmez les tâches cron (toutes les 5 minutes) :
    \f{process_reminders.php}, \f{send_email.php} et
    \f{daily_round_alert.php}, lancées par \texttt{php} en ligne de commande
    (elles refusent les requêtes web).
  \item Connectez-vous avec le compte administrateur initial, puis
    \textbf{changez son mot de passe} et configurez la double
    authentification.
\end{enumerate}
\end{procedure}

\subsection{Mettre à jour l'application}
\begin{procedure}{Mise à jour}
\begin{enumerate}
  \item Sauvegardez la base, \f{storage/} (documents d'éthique et de
    formation) et \f{uploads/} (fichiers des cages).
  \item Récupérez la nouvelle version, puis lancez \texttt{composer
    install}.
  \item Appliquez les migrations : \texttt{php tools/migrate.php}. Utilisez
    un compte administrateur de la base, indiqué par
    \texttt{MIGRATION\_DB\_USERNAME} et \texttt{MIGRATION\_DB\_PASSWORD} dans
    \f{.env}. \texttt{-{}-status} liste les migrations appliquées et en
    attente.
\end{enumerate}
\end{procedure}
\begin{regle}
Une base installée avant l'outil de migrations doit être initialisée une
seule fois en indiquant la dernière migration déjà appliquée, par exemple
\texttt{php tools/migrate.php -{}-baseline=034}. Lancez ensuite
\texttt{php tools/migrate.php}.
\end{regle}

\subsection{Sauvegardes}
Sauvegardez chaque jour la base de données, \f{storage/} et \f{uploads/}, et
conservez les copies hors du serveur. Testez une restauration au moins une
fois par an : les registres réglementaires dépendent de ces sauvegardes.

\section{Paramétrer le laboratoire}
\menu{Paramètres > Gérer le laboratoire}.
\begin{center}
\rowcolors{2}{gray!8}{white}
\begin{tabularx}{\textwidth}{@{}lp{2.5cm}X@{}}
\toprule
\textbf{Réglage} & \textbf{Défaut} & \textbf{Effet}\\
\midrule
Nom, URL & --- & Affichés sur les pages, les étiquettes et dans les e-mails.\\
Fuseau horaire & Europe/Paris & Définit le « jour » de la tournée, des suivis et des rappels.\\
Alerte de la tournée & 18:00 & Heure après laquelle les administrateurs reçoivent la liste des cages non vues.\\
Âge au sevrage & 21 jours & Sert au calcul de la date de sevrage prévue de chaque portée.\\
Rappels du matin & 07:30 & Heure des rappels d'actes expérimentaux et de sevrage.\\
Déconnexion & 30 min & Inactivité tolérée avant fermeture de la session (5 à 720 minutes).\\
Double authentification & obligatoire & Exigée des administrateurs et responsables BEA.\\
Capteurs IoT & --- & Liens des capteurs des salles.\\
\bottomrule
\end{tabularx}
\end{center}
\capture[0.85\linewidth]{a_laboratoire}{Gérer le laboratoire.}

\section{Comptes utilisateurs}
\menu{Paramètres > Gérer les utilisateurs}.
\begin{procedure}{Créer un compte}
\begin{enumerate}
  \item \btn{Créer un utilisateur} : nom complet, e-mail (identifiant),
    fonction, rôle, responsable BEA (oui ou non) et mot de passe initial.
  \item Transmettez le mot de passe initial par un canal distinct de
    l'e-mail, et demandez à la personne de le changer.
  \item Affectez la personne aux cages qu'elle suit : à la création de la
    cage, ou ensuite par le changement de propriété en lot
    (section~\ref{sec:propriete}). Les utilisateurs d'une cage reçoivent ses
    rappels et gèrent ses fichiers.
\end{enumerate}
\end{procedure}
\capture[0.85\linewidth]{a_utilisateurs}{Gestion des utilisateurs.}
\begin{itemize}
  \item \btn{Désactiver} un compte plutôt que le supprimer quand une personne
    part : ses saisies restent attribuées à son nom.
  \item Un compte verrouillé (3 mots de passe ou 5 codes erronés) se
    débloque seul au bout de 15 minutes.
  \item Double authentification perdue :
    \menu{Paramètres > Double authentification}, section administrateur,
    réinitialisation pour l'utilisateur concerné.
\end{itemize}

\section{Propriété des cages}\label{sec:propriete}
\begin{procedure}{Changer le responsable ou les utilisateurs de cages en lot}
\begin{enumerate}
  \item Ouvrez un tableau de bord (\menu{Tableaux de bord > \ldots}) et
    cochez les cages concernées (une, plusieurs, ou toute la page).
  \item Dépliez \btn{Changer la propriété des cages sélectionnées
    (administrateur)}.
  \item Choisissez un \textbf{responsable} (fonction Principal
    Investigator), ou laissez-le inchangé.
  \item Pour les \textbf{utilisateurs}, choisissez l'opération (ajouter,
    remplacer, retirer) et les personnes concernées.
  \item \btn{Changer la propriété}, puis confirmez.
\end{enumerate}
\end{procedure}
\capture[0.9\linewidth]{a_propriete}{Changement de propriété des cages sélectionnées sur un tableau de bord.}
Les cages retirées sont ignorées. Chaque cage modifiée est inscrite au
journal d'activité avec son responsable et ses utilisateurs avant et après.
Lors du départ d'une personne, utilisez \btn{retirer les utilisateurs
choisis} sur toutes ses cages, puis désactivez son compte.

\section{Référentiels}

\subsection{Lignées et souches}
\menu{Paramètres > Gérer les lignées animales} : code et espèce de chaque
lignée. Le code forme le numéro des animaux (\texttt{RSD-0002}). Une espèce
sans lignée ne peut pas être incluse en expérimentation.
\menu{Paramètres > Gérer les souches} : souches et références (RRID).

\subsection{Autorisations de projet}
\menu{Paramètres > Autorisations de projet (n° éthique)}.
\begin{procedure}{Enregistrer une autorisation}
\begin{enumerate}
  \item Saisissez le numéro éthique, le numéro d'autorisation et les dates
    de validité.
  \item Déposez en PDF la demande d'autorisation, l'avis du comité
    d'éthique et la décision d'autorisation. Ils sont stockés hors du web et
    téléchargeables par les seuls utilisateurs connectés.
  \item Indiquez le \textbf{nombre d'animaux autorisés par espèce}.
\end{enumerate}
\end{procedure}
\capture[0.85\linewidth]{a_autorisations}{Autorisations de projet : validité, documents, animaux autorisés par espèce.}
\begin{regle}
La période d'un essai doit être comprise dans la validité de son
autorisation. Aucune inclusion n'est possible si l'autorisation est expirée
ou pas encore valide, ni au-delà des quotas. Les administrateurs et
responsables BEA sont prévenus à 80~\% puis à 100~\% d'un quota.
\end{regle}

\subsection{Grilles d'évaluation clinique}
\menu{Animaux > Grilles d'évaluation clinique} : items et niveaux notés.
Faites valider les grilles par la structure chargée du bien-être animal avant
leur première utilisation. Une grille utilisée devient non modifiable :
dupliquez-la pour la faire évoluer.

\section{Supervision quotidienne}
\begin{procedure}{Routine de l'administrateur}
\begin{enumerate}
  \item Lisez les \textbf{notifications} (cloche) :
    \begin{itemize}
      \item points limites atteints et actions requises ;
      \item signalements sanitaires ;
      \item quotas d'animaux ;
      \item actes expérimentaux et sevrages en retard ;
      \item cages non vues lors de la tournée ;
      \item entrées de formation à valider.
    \end{itemize}
  \item Sur l'\menu{Accueil}, vérifiez les portées à sevrer : celles en
    rouge ont dépassé leur date.
  \item Dans la \menu{Tournée quotidienne}, contrôlez les jours passés pour
    repérer les cages oubliées.
  \item Dans le \menu{Registre des entrées et sorties}, complétez les
    provenances signalées \btn{À compléter}.
\end{enumerate}
\end{procedure}
Les rappels sont produits par \f{process_reminders.php} : la tâche cron doit
tourner. Les e-mails partent par \f{send_email.php}.

\section{Corrections réglementaires}
\begin{trace}
Toutes les corrections ci-dessous demandent un \textbf{motif}. Elles
conservent la valeur d'origine et apparaissent dans les journaux exportés.
Les tables d'audit refusent toute modification ou suppression au niveau de la
base de données.
\end{trace}

\begin{procedure}{Corriger l'entrée ou la sortie d'un animal}
\begin{enumerate}
  \item Ouvrez la fiche de l'animal, cadre \btn{Registre des entrées et
    sorties}.
  \item \btn{Corriger l'entrée (administrateur)} ou \btn{Corriger la sortie
    (administrateur)}. Modifiez les valeurs et saisissez le \btn{Motif de la
    correction}.
  \item Corriger le code de sortie change aussi le statut de l'animal. Les
    inclusions clôturées par la sortie suivent la nouvelle date.
\end{enumerate}
\end{procedure}
\capture[0.85\linewidth]{a_registre}{Registre des entrées et sorties.}
Les corrections apparaissent sur la fiche et dans le dossier BEA. L'animal
est marqué « corrigé » dans le registre.

\begin{itemize}
  \item \textbf{Suivis journaliers} et \textbf{confirmations d'actes} : se
    corrigent depuis le suivi, avec un motif. Les valeurs remplacées sont
    conservées. Pour un acte, le médicament est recrédité.
  \item \textbf{Administrations de médicaments} : annulation unique avec
    motif. Une ordonnance ne se modifie jamais ; il faut en créer une
    nouvelle.
  \item \textbf{Suppression définitive d'un animal} : réservée aux erreurs de
    saisie manifestes (animal créé par erreur). Elle demande de retaper le
    numéro de l'animal. Elle est journalisée avant la suppression. Préférez
    toujours une sortie avec son code.
\end{itemize}

\section{Volet BEA}
\begin{itemize}
  \item \menu{Paramètres > Analyse rétrospective BEA} : parcours par essai,
    plan, puis expérience. Les sévérités réelles (sans réveil, légère,
    modérée, sévère) sont classées en regard de la sévérité attendue du
    plan. Le dossier complet de chaque animal est imprimable.
  \item \menu{Paramètres > Bilan annuel d'utilisation des animaux} :
    effectifs, procédures terminées par sévérité réelle, utilisation par
    autorisation, sorties, animaux non utilisés. Export CSV ou ZIP.
  \item \menu{Paramètres > Livrets de formation du personnel} : validez ou
    refusez, avec un motif, les entrées déclarées. Réglez l'obligation de
    formation continue (3 jours sur 6 ans par défaut). Suivez les personnes
    à traiter en priorité. Une personne non encore compétente ne peut
    confirmer un acte qu'en indiquant son encadrant.
\end{itemize}
\capture[0.85\linewidth]{a_retrospective}{Analyse rétrospective BEA d’une expérience.}

\section{Pharmacie et stupéfiants}
\begin{itemize}
  \item \btn{Nouveau médicament} : cochez \btn{Stupéfiant} pour un produit
    soumis au registre spécial, et indiquez sa concentration (par exemple
    100~mg/mL) pour le calcul des doses selon le poids.
  \item \btn{Réception d'un lot} : numéro, péremption, quantité,
    emplacement. \btn{Ajuster} un stock exige un motif.
  \item \btn{Registre des stupéfiants} : tous les mouvements par lot et leur
    solde. Les administrateurs y enregistrent les \textbf{inventaires} :
    quantité comptée, écart justifié, appliqué au stock. Un écart entre le
    registre et le stock est signalé.
\end{itemize}
\capture[0.85\linewidth]{a_stupefiants}{Registre des stupéfiants.}

\section{Exports, audit et inspection}
\begin{itemize}
  \item \menu{Paramètres > Export brut des journaux} : journal d'activité,
    événements et mouvements des animaux, contrôles et cycle de vie des
    cages, portées, suivi sanitaire, actions de point limite, actes et leurs
    corrections, corrections des suivis et du registre, formations,
    terrariums, notifications, e-mails. Export CSV par journal ou ZIP, avec
    filtre de dates. Chaque export est tracé.
  \item \menu{Paramètres > Exporter en CSV} : toutes les tables (sans les
    mots de passe).
  \item \menu{Paramètres > Export pour la V2} : export JSON sans secrets
    (mots de passe, jetons, double authentification, corps des e-mails).
\end{itemize}
\capture[0.85\linewidth]{a_journaux}{Export brut des journaux.}

\begin{procedure}{Préparer une inspection}
\begin{enumerate}
  \item \menu{Registre des entrées et sorties} de la période : export PDF ou
    CSV, provenances complètes, corrections justifiées.
  \item \menu{Registre sanitaire} et \btn{Registre des stupéfiants} avec le
    dernier inventaire.
  \item Pour chaque projet : autorisation et documents PDF, quotas
    consommés, rapports d'expérience, dossiers animaux.
  \item Bilan annuel et analyse rétrospective BEA.
  \item Livrets de formation à jour (synthèse du personnel).
\end{enumerate}
\end{procedure}

\section{Sécurité et exploitation}
\begin{itemize}
  \item Ne désactivez pas la double authentification des comptes
    privilégiés.
  \item Ne partagez jamais un compte : la traçabilité repose sur l'auteur de
    chaque saisie.
  \item Après toute intervention sur le serveur, vérifiez que \f{/.env}
    répond 403 ou 404.
  \item La CI du dépôt (audit des dépendances, PHPStan, tests) doit être
    verte avant chaque mise en production.
  \item En cas de besoin, le journal d'activité permet de retrouver qui a
    fait quoi et quand.
\end{itemize}

\end{document}
